\textbf{Problème 5.3 (Dujella, 2026).} \begin{enumerate}
\item Existe-t-il un quadruplet diophantien polynomial ayant la propriété $D((4k-1)(4k+1))$, en dehors des quadruplets obtenus à partir des deux quadruplets diophantiens polynomiaux connus ayant la propriété $D(4k+3)$ ?
\item Soit $n=(4k-1)(4k+1)$, où $k>3$ est un entier positif tel que $4k-1$, $4k+1$, $(n-1)/2$, $(n-9)/2$, $(9n-1)/2$, $16k-5$, $16k+5$ soient tous premiers (des exemples de tels $n$ sont $9294986317823$, $27363653936783$, $48073810526783$). Existe-t-il un tel $n$ pour lequel il existe plus de deux $D(n)$-quadruplets ?
\end{enumerate}

\medskip

\textit{Remarque.} Si l'un des nombres $4k-1$, $4k+1$, $(n-1)/2$, $(n-9)/2$, $(9n-1)/2$ n'est pas premier, alors des $D(n)$-quadruplets supplémentaires (à éléments positifs) peuvent être obtenus d'après Dujella, A. : A problem of Diophantus and Dickson's conjecture. In : Győry, K., Pethő, A., Sós, V.T. (éds.) Number Theory, Diophantine, Computational and Algebraic Aspects, pp. 147--156. Walter de Gruyter, Berlin (1998) ; si $16k-5$ ou $16k+5$ n'est pas premier, il existe des $D(n)$-quadruplets supplémentaires à signes mêlés (Dujella (2026)).
